\begin{center}
\textbf{\large{Abstract}}
\end{center}
\addcontentsline{toc}{chapter}{Abstract}
\begin{small}

Database management systems (DBMS) are critical infrastructure components, and their vulnerabilities represent serious security threats. Automated vulnerability detection through fuzzing is a promising approach, but existing grammar-based fuzzers apply uniform random sampling over production rules, wasting mutation budget on low-value syntactic paths and failing to prioritize the structural patterns most likely to expose bugs.

This thesis presents a grammar-based greybox fuzzing system targeting SQLite. A structural primitives grammar of 475 rules encodes SQL patterns known to trigger vulnerability classes---including integer overflow in expression trees, use-after-free in query planner optimizations, and memory corruption in aggregate processing---without hardcoding specific proof-of-concept inputs. On top of this grammar, a multi-armed bandit controller using the EXP3 algorithm adaptively selects grammar rules based on coverage feedback from an AFL-compatible fork server, aiming to concentrate mutations on rules that generate coverage-increasing inputs.

Experiments across 53 campaigns on 3 SQLite versions containing 6 known CVEs show that uniform sampling discovers significantly more unique crash classes than the bandit controller (mean 8.6 vs.\ 4.2 per campaign, $p \approx 0.01$, Mann-Whitney U test). Analysis reveals that the bandit converges to exploitation of a small rule subset, starving exploration and reducing grammatical diversity. A grammar gap analysis further shows that 4 of the 6 target CVEs were unreachable until structural primitives for window functions and self-referential generated columns were added in grammar version 3.2, demonstrating that grammar coverage is the dominant factor for vulnerability reachability.

\vspace*{1cm}
\textbf{Keywords:} Grammar-based Fuzzing, Greybox Fuzzing, DBMS Security, Vulnerability Detection, Multi-Armed Bandit
\end{small}
